\documentclass[9pt]{beamer}
\usepackage{luatextra}

%% Modules nécessaires pour utiliser le thème ecl24
\usepackage{tikz}
\usepackage{ifthen,etoolbox}
\graphicspath{{IMAGES/}}

%% Pour utiliser le thème ecl24
\usetheme[ccbyncnd]{ecl24}

%% Afin de supprimer les symbols de navigation Beamer
\beamertemplatenavigationsymbolsempty

\begin{document}

%% Informations utilies pour la page de garde et les pieds de page
\title{Thème Beamer Latex de l'école centrale de Lyon}
\subtitle{Charte graphique de l'école centrale de Lyon 2024}
\author[Dr Bichot]{Dr Charles-Edmond Bichot} % shortauthor est utilisé pour ccbyncnd
\institute{École centrale de Lyon}
\date{22 janvier 2026}

%% Création de la page de garde
\maketitle

\section{Le thème \emph{ecl24}}

\begin{frame}[fragile]{\insertsection} % 'fragile' nécessaire à verbatim
  %% Bloc important
  \begin{alertblock}{Prérequis}
    Pour utiliser le thème Beamer \emph{ecl24}, il vous faut appeler les modules \emph{tikz}, \emph{ifthen} et \emph{etoolbox}.

    Pensez à vérifier que les images utilisée par le thème Beamer \emph{ecl24} sont accessibles en indiquant par exemple le chemin vers ces images.
  \end{alertblock}

  %% Bloc d'exemple
  \begin{exampleblock}{Utilisation du thème \emph{ecl24}}
    Ces commandes sont à insérer dans l'en-tête de votre fichier latex (\emph{ie} avant \verb+\begin{document}+) :
    \begin{verbatim}
      \usepackage{tikz}     
      \usepackage{ifthen,etoolbox}
      \graphicspath{{IMAGES/}}
      \usetheme{ecl24}
    \end{verbatim}
  \end{exampleblock}
\end{frame}


\section{Licence creative commons BY-NC-ND}

\begin{frame}[fragile]{\insertsection}
  %% Bloc normal
  \begin{block}{Utiliser la licence creative commons BY-NC-ND}
    Pour utiliser la licence creative commons BY-NC-ND, activez l'option \emph{ccbyncnd} en appelant le thème \emph{ecl24} :

    \verb+\usetheme[ccbyncnd]{ecl24}+
  \end{block}
\end{frame}

\end{document}
